\documentclass[11pt]{article}
\usepackage[T1]{fontenc}
\usepackage{lmodern,amsmath,amssymb,amsthm,mathtools}
\usepackage[margin=1in]{geometry}
\usepackage{microtype}
\usepackage[numbers,sort&compress]{natbib}
\usepackage{xurl}
\usepackage[hidelinks]{hyperref}
\hypersetup{pdftitle={Compact short-time support for fractional-pressure diffusion},pdfauthor={Kaiwen Jin}}
\newcommand{\R}{\mathbb R}
\newcommand{\dd}{\,\mathrm d}
\newcommand{\eps}{\varepsilon}
\newtheorem{theorem}{Theorem}[section]
\newtheorem{proposition}[theorem]{Proposition}
\newtheorem{lemma}[theorem]{Lemma}
\theoremstyle{definition}
\newtheorem{definition}[theorem]{Definition}
\numberwithin{equation}{section}
\title{An energy weak solution with compact short-time support\protect\\
 for fractional-pressure diffusion}
\author{Kaiwen Jin\\\small\texttt{yui.ui.siki@gmail.com}}
\date{November 2026}

\begin{document}
\maketitle
\begin{abstract}
We construct a global bounded energy weak solution of
$u_t=\nabla\cdot(u^{m-1}\nabla(-\Delta)^{-s}u)$ in three dimensions,
with $m=3/2$ and one $s\in(7/8,1)$, whose support remains compact for a
positive time. The solution conserves mass, obeys the initial supremum
bound, is continuous into weak $L^1$, and satisfies the pressure energy
and every dissipative $L^q$ inequality from the initial time. Thus the
energy class specified in AIM Question~468 does not force instantaneous
spreading. The construction retains a moving jump to vacuum. Its positive
boundary entropy contribution is compensated by an interior compressive
shock at higher densities.
A weighted interaction estimate separates the nonlocal entropy production
from the viscosity measure and identifies enough dissipation at that shock
to dominate the boundary contributions for all powers on a single time interval.
\end{abstract}

\section{Introduction and main result}\label{sec:intro}

Nonlocal pressure makes the support of a diffusing density sensitive to
both the mobility and the solution concept. We consider
\begin{equation}\label{eq:model}
 u_t=\nabla\cdot(u^{m-1}\nabla p),\qquad
 p=(-\Delta)^{-s}u,\qquad x\in\R^N,
\end{equation}
where $m>1$, $0<s<1$, and the inverse fractional Laplacian has Fourier
multiplier $|\xi|^{-2s}$. This is the porous-medium model with fractional
potential pressure studied by Stan, del Teso and V\'azquez
\cite{Stan2015,Stan2016,Stan2018}. The pressure couples distant regions,
while the mobility degenerates at vacuum; deciding whether a compact
front persists is therefore a qualitative question beyond global existence.

The global existence theory supplies mass and supremum bounds and two
families of energy estimates \cite[Theorem~2.2]{Stan2016}. Propagation
depends on the parameter range. The survey \cite[Theorems~5.3--5.4]{Stan2018}
records finite propagation for constructed solutions when $m\geq2$, and
infinite propagation in dimension one when $1<m<2$. The latter uses an
integrated scalar problem; Remark~7 there identifies the multidimensional
extension as an open issue. A result for one selected construction and a
statement about every solution in a given energy class are distinct
questions. The formulation in AIM Question~468 \cite{AIM468} explicitly
asks the second question for $N\geq2$ and $1<m<2$.

We state its solution class precisely. In the following definition the
initial datum is nonzero, nonnegative and compactly supported.

\begin{definition}\label{def:class}
Let $N\geq2$, $1<m<2$, $0<s<1$, and
$d\in L^1(\R^N)\cap L^\infty(\R^N)$.
A bounded energy weak solution from $d$ is a nonnegative distributional
solution of \eqref{eq:model} with locally integrable flux, such that
\[
 t\longmapsto\int_{\R^N}u(t,x)\psi(x)\,\dd x
 \quad\hbox{is continuous for every }\psi\in L^\infty(\R^N),
 \qquad u(0)=d,
\]
and, at every $t\geq0$, $\int u(t)=\int d$ and
$\|u(t)\|_\infty\leq\|d\|_\infty$. For every $t>0$ it satisfies
\begin{equation}\label{eq:class-pressure}
 \frac12\|(-\Delta)^{-s/2}u(t)\|_2^2+
 \int_0^t\int u^{m-1}|\nabla(-\Delta)^{-s}u|^2
 \leq\frac12\|(-\Delta)^{-s/2}d\|_2^2,
\end{equation}
and, for every real $q>1$ and $t>0$,
\begin{equation}\label{eq:class-power}
 \|u(t)\|_q^q+\frac{4q(q-1)}{(m+q-1)^2}
 \int_0^t\|(-\Delta)^{(1-s)/2}u^{(m+q-1)/2}\|_2^2\,\dd r
 \leq\|d\|_q^q.
\end{equation}
\end{definition}

These are inequalities from the initial time. Local entropy admissibility
or membership in a specified approximation limit would impose additional
conditions on this definition. The question is whether every solution in
this class has
$\int_{|x|>R}u(t,x)\dd x>0$ for every finite $R>0$ and every $t>0$.
This expresses unbounded essential support and does not require pointwise
positivity. Our result gives a negative answer to that universal assertion.

\begin{theorem}\label{thm:main}
There exist $0<\alpha<1/4$, $s=1-\alpha/2$, $T_*>0$, and a global
bounded energy weak solution on $\R^3$ with $m=3/2$ and initial datum
$F$ given by \eqref{eq:datum}, such that $\|F\|_\infty=100$ and
\[
 u(t,x)=0\quad\hbox{for almost every }|x|>R(t),
 \qquad 0<t<T_*.
\]
Here $R$ is a positive finite $C^1$ function on $[0,T_*]$ with
$R(0)=101/100$. The solution is radial on this initial interval.
In particular the universal instantaneous-spreading assertion in
Question~468 is false for the class of Definition~\ref{def:class}.
\end{theorem}

All the inequalities in the definition hold for every terminal time, with
one time interval supporting all exponents. The construction gives one
choice of parameters and solution; it makes no propagation assertion
about every other solution or about the whole range $1<m<2$, $0<s<1$.
This quantifier distinction is enough to answer the stated universal
question, and leaves the behavior of separately selected solution classes
as a different problem.

The reason a compact front can meet these global inequalities is visible
in its entropy balance. A downward jump from a positive density to vacuum
creates a positive boundary entropy contribution. We put an upward jump at a smaller
radius, so the concave transport flux creates a compressive scalar shock.
The internal shock dissipates at densities strictly above those entering
the boundary contribution. One weighted dissipation bound can consequently dominate
the boundary term for every power, including exponents for which the fractional inequality
has no strict remainder.

Two analytic points are needed to make this comparison rigorous. The
moving viscous ball must exist on a time interval independent of its
viscosity; a boundary layer adapted to the singular fractional source
provides the required normal-flux control. At the internal jump the
nonlocal production must be distinguished from surviving viscosity.
An interaction estimate gives fourth-power spatial increments, which
make the nonlocal defect strongly compact in $L^2$. At the hyperbolic
scale of the initial jump this defect vanishes, and scalar entropy
uniqueness identifies the viscosity production with that of a single
compressive shock. This is the comparison used to prove the exact
inequalities in Definition~\ref{def:class}.

Section~\ref{sec:construction} constructs the common-time viscous family
and its compact limit. Section~\ref{sec:interaction} proves the
interaction and defect estimates. Section~\ref{sec:shock} identifies the
initial shock and proves all power inequalities. Section~\ref{sec:continuation}
extends the solution globally and completes the theorem.
We use the unitary Fourier transform
$\widehat w(\xi)=(2\pi)^{-N/2}\int e^{-ix\cdot\xi}w(x)\dd x$.
Sobolev spaces on bounded domains are restriction spaces, and
$H^{-1}=(H^1_0)^*$. Constants denoted by $K$ may change between estimates;
after its selection, $\alpha$ is fixed throughout the construction.

\section{Construction on a moving ball}\label{sec:construction}

We fix
\[
 a=\frac12,\qquad A=\frac1{16},\qquad U=100,\qquad
 S=1,\qquad R_0=\frac{101}{100}.
\]
The small positive order $\alpha$ will be chosen below. We write
$s=1-\alpha/2$, $\beta=1-\alpha$, and
\[
 L_\alpha=(-\Delta)^{\alpha/2},\qquad
 K_\alpha=(-\Delta)^{-1+\alpha/2},\qquad
 v[w]=-\nabla K_\alpha w,\qquad f(k)=k^a.
\]
Thus the equation is $u_t+\nabla\cdot(v[u]f(u))=0$ and
$\nabla\cdot v[u]=L_\alpha u$. Densities defined on a ball are always
extended by zero when a whole-space operator is applied.

\subsection{The profile and the pressure field}

Set $h(z)=10z^3-15z^4+6z^5$ on $[0,1]$, and define the radial datum
\begin{equation}\label{eq:datum}
 F(r)=\begin{cases}
 1,&0\leq r<S,\\
 U-(U-A)h\bigl((r-S)/(R_0-S)\bigr),&S<r<R_0,\\
 0,&r>R_0.
 \end{cases}
\end{equation}
Its values on the two spheres are immaterial. There is an upward jump from
$1$ to $U$ at $S$, and a downward jump from $A$ to zero at $R_0$.
The first two derivatives of the annular profile vanish at both endpoints.
In particular $F$ is nonzero, nonnegative, compactly supported and in
$BV\cap L^1\cap L^\infty$, with $\|F\|_\infty=U$.
Every positive power of $F$ belongs to $H^{\alpha/2}(\R^3)$ for $0<\alpha<1$.
Indeed it is piecewise Lipschitz, and flattening a jump sphere bounds the
cross-interface part of its Gagliardo integral by a constant times
\[
 \int_0^1\int_0^1(d_1+d_2)^{-1-\alpha}\,\dd d_1\dd d_2<\infty.
\]
The same-region part follows from Lipschitz continuity, and the part away
from the diagonal follows from bounded support.

For $0\leq\alpha<1$ the normalized potential and velocity are
\begin{align}
 K_\alpha w(x)&=d_\alpha\int_{\R^3}|x-y|^{-1-\alpha}w(y)\,\dd y,
 &d_\alpha&=\frac{\Gamma((1+\alpha)/2)}
 {2^{2-\alpha}\pi^{3/2}\Gamma(1-\alpha/2)},\label{eq:riesz}\\
 v[w](x)&=c_\alpha\int_{\R^3}
 \frac{x-y}{|x-y|^{3+\alpha}}w(y)\,\dd y,
 &c_\alpha&=(1+\alpha)d_\alpha.\label{eq:velocity}
\end{align}
These integrals are continuous for bounded integrable densities.
The velocity kernel is locally integrable since
$\int_0^1r^{-\alpha}\,\dd r<\infty$. Its near-field contribution on a ball
of radius $\delta$ is at most $K\|w\|_\infty\delta^{1-\alpha}$;
the corresponding potential contribution is at most
$K\|w\|_\infty\delta^{2-\alpha}$. The far kernels are bounded, so mass
controls the remaining contributions. For fixed $0<\alpha<1$, splitting a
velocity difference at twice the displacement gives
\begin{equation}\label{eq:velocity-holder}
 |v[w](x)-v[w](x')|\leq K\|w\|_\infty |x-x'|^{1-\alpha}
\end{equation}
on a common bounded support, with the far-field contribution controlled by
mass. Constants in this estimate need not stay bounded as $\alpha\downarrow0$.

For clarity, the Fourier normalization in \eqref{eq:riesz} can be checked
without an operator convention left implicit. Insert
\[
 |z|^{-1-\alpha}=\frac1{\Gamma((1+\alpha)/2)}
 \int_0^\infty t^{(1+\alpha)/2-1}e^{-t|z|^2}\,\dd t.
\]
The unitary Fourier transform of the Gaussian is
$(2t)^{-3/2}e^{-|\xi|^2/(4t)}$. Substitution of $|\xi|^2/(4t)$,
including the convolution factor $(2\pi)^{3/2}$, gives the multiplier
$|\xi|^{\alpha-2}$. Testing first on Schwartz functions and then truncating
the locally integrable kernels justifies the identity in distributions.
Differentiation gives \eqref{eq:velocity} and
\begin{equation}\label{eq:divv}
 -\Delta K_\alpha w=\nabla\cdot v[w]=L_\alpha w.
\end{equation}
The singular-integral normalization, in dimension $n=1$ or $3$, is
\begin{equation}\label{eq:frac-kernel}
 L_{n,\alpha}w(x)=c_{n,\alpha}\,\operatorname{PV}
 \int\frac{w(x)-w(y)}{|x-y|^{n+\alpha}}\,\dd y,
 \qquad c_{n,\alpha}=
 \frac{\alpha2^{\alpha-1}\Gamma((n+\alpha)/2)}
 {\pi^{n/2}\Gamma(1-\alpha/2)}.
\end{equation}
One obtains this coefficient by the same Gaussian representation applied to
$1-\cos(z\cdot\xi)$ and one integration by parts. Symmetric truncation and
Plancherel identify the positive quadratic form with
$\|(-\Delta)^{\alpha/4}w\|_2^2$.

Two further bounds will be used repeatedly. Inside a ball $D$,
\begin{equation}\label{eq:indicator}
 0\leq L_\alpha\mathbf1_D(x)\leq K\operatorname{dist}(x,\partial D)^{-\alpha}.
\end{equation}
Also $L_\alpha$ maps bounded compactly supported $C^\theta$ functions to
bounded functions if $\theta>\alpha$. Both statements follow by splitting
the singular integral at the distance to the boundary or at one. The
radial integrals are respectively $\int_d^\infty r^{-1-\alpha}\dd r$ and
$\int_0^1r^{\theta-\alpha-1}\dd r$.

For radial $w$ write $v[w](r)$ for the outward component, and put
\[
 W[w;r]=4\pi r^2v[w](r).
\]
If $w\geq0$ is supported in $B_r$ and has mass $M>0$, then at $x=rn$,
$|n|=1$, the inequality $r-y\cdot n\geq |rn-y|^2/(2r)$ gives
\begin{equation}\label{eq:boundary-positive}
 v[w](r)\geq c_\alpha M(2r)^{-2-\alpha}>0.
\end{equation}
At $\alpha=0$, Gauss' formula yields
$v_0[w](r)=r^{-2}\int_0^r\rho^2w(\rho)\,\dd\rho$.
Consequently
\[
 W_0[F;S]=\frac{4\pi}3,\qquad W_0[F;R_0]=\int F,
 \qquad
 \frac{W_0[F;R_0]}{W_0[F;S]}
 \leq1+100\bigl((101/100)^3-1\bigr)=4.0301.
\]
For this fixed datum the two velocity values converge to their Newtonian
values as $\alpha\downarrow0$. A majorant with any fixed order less than
one controls the near kernel, and dominated convergence applies on the
remaining compact support. Thus both weights are positive and their ratio
is less than five for all sufficiently small positive $\alpha$.

The concavity deficits that will be compared are
\begin{align}
 g_A(k)&=k^a-A^{a-1}k, &0\leq k\leq A,\\
 H_U(k)&=k^a-1-\frac{U^a-1}{U-1}(k-1), &1\leq k\leq U.
\end{align}
They are positive in the interiors by strict concavity. Direct integration gives
\begin{equation}\label{eq:margin}
 G_A:=\int_0^A\frac{g_A(k)}k\,\dd k=\frac14,
 \qquad J:=\int_1^U\frac{H_U(k)}k\,\dd k
 =9-\frac{10}{11}\log100>4.
\end{equation}
We henceforth fix one $0<\alpha<1/4$ sufficiently small that
\begin{equation}\label{eq:strict-margin}
 W[F;R_0]G_A<W[F;S]J.
\end{equation}
There is no numerical endpoint asserted for this choice. Notice also that,
for every $q>1$,
\begin{equation}\label{eq:density-cost}
 \int_0^A k^{q-2}g_A(k)\,\dd k\leq A^{q-1}G_A.
\end{equation}

\subsection{A viscous free boundary problem}

We approximate only the internal jump of $F$. There are radial
$g_j\in C^2(\overline B_{R_0})$, $A\leq g_j\leq U$, with
$g_j=A$, $\partial_ng_j=0$ on the boundary, whose zero extensions converge
to $F$ in every finite $L^q$ and in $L^1$. They agree with $F$ in one fixed
outer collar, with a common $C^1$ bound there. To construct them, mollify
the radial profile near $S$ and join it to $F$ with a fixed smooth cutoff.
Only values in $[A,U]$ are averaged. Their masses lie in a common interval
$[M_-,M_+]$ with $M_->0$, after deleting a finite prefix if necessary.
Choose $\eps_j=2^{-j}$, with a harmless relabeling after that deletion.

For fixed $\eps>0$ and one such datum consider, in $D(t)=B_{R(t)}$,
\begin{align}
 w_t&=\nabla\cdot(w^a\nabla K_\alpha w)+\eps\Delta w,
 &w|_{\partial D(t)}&=A,\label{eq:viscous}\\
 R'&=A^{a-1}v[w](R)-\frac\eps A\partial_nw,
 &R(0)&=R_0,\qquad w(0)=g.\label{eq:stefan}
\end{align}
The positive boundary value permits a discontinuity to exterior vacuum.
The law \eqref{eq:stefan} is precisely the cancellation required by mass
conservation for the zero extension.

We first give the local realization used in subsequent integrations by
parts. Choose $3<p_0<1/\alpha$. On $L^{p_0}(B_1)$ let
\[
 Gz=\Delta z,\qquad D(G)=W^{2,p_0}(B_1)\cap W^{1,p_0}_0(B_1),
 \qquad \Sigma(t)=e^{tG}.
\]
The Dirichlet Laplacian is densely defined and sectorial
\cite[Section~2.4]{Lunardi2004}. Its analytic semigroup satisfies
\begin{align}
 \|\Sigma(t)h\|_{C^1}&\leq K(1+t^{-\lambda})\|h\|_{p_0},
 &\lambda&=\frac12+\frac3{2p_0}<1,\label{eq:heat-smoothing}\\
 \|[\Sigma(h)-I]\Sigma(r)\|_{L^{p_0}\to C^1}
 &\leq K\min(r^{-\lambda},hr^{-1-\lambda}).\label{eq:heat-increment}
\end{align}
Here and below bounded time intervals are understood in heat estimates.
The semigroup is uniformly bounded and strongly continuous on the space of
$C^1$ functions with zero boundary value, in the full $C^1$ norm.
This space imposes no zero normal derivative.

Here are details of the boundary and $C^1$ assertions. Comparison with the
whole-space Gaussian gives the $L^{p_0}$ to $L^\infty$ estimate. At a tangent
plane to the convex ball, comparison with
$\|h\|_\infty\operatorname{erf}(\ell/(2\sqrt t))$, where $\ell$ is the
inward affine distance, bounds the boundary gradient by
$\|h\|_\infty/\sqrt{\pi t}$. For a heat solution,
\[
 (\partial_t-\Delta)(t|\nabla z|^2+z^2)
 =-|\nabla z|^2-2t|D^2z|^2.
\]
The maximum principle and semigroup factorization give
\eqref{eq:heat-smoothing}. For zero-trace $C^1$ data with Lipschitz constant
$L$, the stationary barriers $\pm L\ell$ instead give a uniform gradient
bound. For smooth such data,
$\Sigma(t)h-h=\int_0^t\Sigma(r)\Delta h\,\dd r$ tends to zero in $C^1$
by $\lambda<1$. Smooth zero-trace approximation and the uniform bound
extend this to all zero-trace $C^1$ data. Finally the analytic estimate
$\|G\Sigma(r)\|_{L^{p_0}\to L^{p_0}}\leq K/r$
\cite[Theorem~1.2.3]{Lunardi2004}, combined with
\eqref{eq:heat-smoothing}, proves \eqref{eq:heat-increment} by integration.

To solve \eqref{eq:viscous}--\eqref{eq:stefan}, put
$y=x/R(t)$, $\dd\tau/\dd t=\eps/R(t)^2$, $\rho(\tau)=R(t)$ and
$z(\tau,y)=w(t,x)$. Write
$\theta=\rho^{2-\alpha}/\eps$ and
$\widetilde v=v[z\mathbf1_{B_1}]$. The equations become
\begin{align}
 z- A&=\Sigma(\tau)(g(R_0\cdot)-A)
       +\int_0^\tau\Sigma(\tau-r)\mathcal N(z(r),\rho(r))\,\dd r,
                                                        \label{eq:mild}\\
 \rho_\tau&=\theta\rho A^{a-1}\widetilde v(1)
              -\frac\rho A\partial_nz(1),\label{eq:rho}\\
 \mathcal N(z,\rho)&=
 \left[\theta A^{a-1}\widetilde v(1)y
 -\frac{\partial_nz(1)}A y-a\theta z^{a-1}\widetilde v(y)\right]\cdot\nabla z
 \nonumber\\[-2mm]
 &\hspace{8mm}-\theta z^a
 \left(L_\alpha[(z-A)\mathbf1_{B_1}]+A L_\alpha\mathbf1_{B_1}\right).
                                                        \label{eq:forcing}
\end{align}
For positive radial $C^1$ functions with trace $A$ and positive radii,
$\mathcal N$ is locally Lipschitz into $L^{p_0}$. Indeed the zero extension
of $z-A$ is globally Lipschitz, so its fractional Laplacian is bounded;
\eqref{eq:indicator} is in $L^{p_0}$ because $\alpha p_0<1$.
All boundary derivatives and the remaining products are locally Lipschitz
in the indicated $C^1$ norm. The radius vector field has the same property.
On continuous $C^1$ density paths and continuous radius paths, the two
integral equations have contraction constants bounded by
$K\tau_0^{1-\lambda}$ and $K\tau_0$. Strong $C^1$ continuity of the linear
path permits a positive closed ball of paths on a sufficiently short
interval. The contraction gives a unique positive radial local solution.

This mild solution has the required strict realization. The increment
estimate \eqref{eq:heat-increment}, split at $r=h$, gives a time
$C^{1-\lambda}$ bound in $C^1$ for the Duhamel term. The initial datum
$g(R_0\cdot)-A$ belongs to $D(G)$, and the same bound follows for its heat
path by integrating $\Sigma(r)G(g(R_0\cdot)-A)$. Hence
$\mathcal N(z,\rho)$ is H\"older continuous in time into $L^{p_0}$.
The strict-solution theorem \cite[Theorem~4.2.6(ii)]{Lunardi2004} applies:
the compatibility condition is
$G(g(R_0\cdot)-A)+\mathcal N(z(0),R_0)\in\overline{D(G)}$,
which is automatic since $D(G)$ is dense in $L^{p_0}$. Thus
\[
 z-A\in C^1_\tau L^{p_0}\cap C_\tau D(G).
\]
Interior localization and the Gaussian heat formula give classical
$C^1$ time and $C^2$ space regularity for positive times: after subtracting
the forcing value at the evaluation point, the differentiated heat kernel
has an integrable time singularity by the positive H\"older exponent.
The established $W^{2,p_0}$ realization gives spatial
$C^{1,1-3/p_0}$ regularity by Sobolev embedding. Thus the drift--gradient
product is H\"older on interior subballs, and the fractional kernel
estimates also make the remaining forcing terms spatially H\"older.
This also treats the origin in Cartesian
coordinates. Finally
\[
 v[w](Ry)=R^{1-\alpha}\widetilde v(y),\qquad
 L_\alpha w(Ry)=R^{-\alpha}L_\alpha[z\mathbf1_{B_1}](y),
 \qquad \frac{\dd t}{\dd\tau}=\frac{\rho^2}\eps
\]
undo the transformation and give exactly
\eqref{eq:viscous}--\eqref{eq:stefan}. Regularity here concerns the interior
and its one-sided boundary traces; the exterior jump is retained.

\subsection{Identities at the moving boundary}

Put $j=w^a\nabla K_\alpha w+\eps\nabla w$ inside $D(t)$.
The boundary law says $AR'=-j\cdot n$. Reynolds' rule therefore gives
\begin{equation}\label{eq:extended-mass}
 \partial_tw=\nabla\cdot(j\mathbf1_{D(t)}),\qquad
 \int w(t)=\int g
\end{equation}
for the zero extension, including its initial trace. At an interior global
maximum, $L_\alpha w\geq0$, $\nabla w=0$ and $\Delta w\leq0$.
The equation and a strict first-contact argument imply $w\leq U$ as long
as the positive local solution exists.

For an entropy $\eta$ with $\eta(0)=0$, define
\begin{equation}\label{eq:entropy-transforms}
 \Psi_\eta(k)=\int_0^k\eta''(z)f(z)\,\dd z,
 \quad F_\eta(k)=\int_0^k\Psi_\eta(z)\,\dd z,
 \quad Q_\eta(k)=\eta'(k)f(k)-\Psi_\eta(k).
\end{equation}
For local distributional identities we normalize also $\eta'(0)=0$.
The powers $\eta(k)=k^q$, $q>1$, and $k\log k$ are admitted by
approximation whenever the displayed integrals at zero converge.
For increasing $G$ set
\begin{equation}\label{eq:pair}
 \mathcal B_\alpha(G,w)=\frac{c_{3,\alpha}}2
 \iint\frac{(G(w(x))-G(w(y)))(w(x)-w(y))}
 {|x-y|^{3+\alpha}}\,\dd x\dd y.
\end{equation}
This is a nonnegative extended integral. When finite it is the symmetric
form of $\int G(w)L_\alpha w$.

Reynolds' rule and integration by parts yield the exact identity
\begin{align}
 \frac\dd{\dd t}\int_D\eta(w)
 +\mathcal B_\alpha(\Psi_\eta,w)
 +\eps\int_D\eta''(w)|\nabla w|^2
 ={}&W[w;R]\int_0^A\eta''(k)g_A(k)\,\dd k\nonumber\\
 &+\eps|\partial B_R|\partial_nw
       \left(\eta'(A)-\frac{\eta(A)}A\right).
                                                        \label{eq:entropy-global}
\end{align}
Here $W[w;R]=4\pi R^2v[w](R)$ and
$g_A(k)=k^a-A^{a-1}k$.
In fact the initial boundary term is
$j\cdot n[\eta'(A)-\eta(A)/A]$. Since
$\nabla\Psi_\eta(w)=\eta''(w)w^a\nabla w$ and
$\nabla\cdot\nabla K_\alpha w=-L_\alpha w$, integrating the pressure
term contributes $W[w;R]\Psi_\eta(A)$.
The remaining coefficient is
\[
 \Psi_\eta(A)-A^a\eta'(A)+A^{a-1}\eta(A)
 =\int_0^A\eta''(k)g_A(k)\,\dd k.
\]
This proves \eqref{eq:entropy-global}, including its sign. For fixed
$\eps$ the established strict $L^{p_0}$ realization, positive density and
$C^1$ traces justify the time-integrated calculation; the boundary
singularity has size $d^{-\alpha}$ and is integrable. Entropy approximation
on positive interior values gives the asserted non-smooth-at-zero cases.
For powers,
\begin{equation}\label{eq:psi-power}
 \Psi_{k^q}(k)=\frac{q(q-1)}{a+q-1}k^{a+q-1}.
\end{equation}
For the logarithmic entropy, \eqref{eq:entropy-global} reads
\begin{equation}\label{eq:log-identity}
 \frac\dd{\dd t}\int w\log w+
 \frac1a\int w^aL_\alpha w+
 \eps\int_D\frac{|\nabla w|^2}{w}
 =W[w;R]G_A+\eps|\partial B_R|\partial_nw.
\end{equation}

The pressure energy is
\[
 \mathcal E_\alpha(w)=\frac12\|(-\Delta)^{-s/2}w\|_2^2
 =\frac12\int wK_\alpha w.
\]
It is finite for bounded compact densities, since the potential kernel is
locally integrable and bounded away from its singularity. Its identity is
\begin{equation}\label{eq:pressure-identity}
 \frac\dd{\dd t}\mathcal E_\alpha(w)
 +\int_Dw^a|\nabla K_\alpha w|^2
 +\eps\int wL_\alpha w=\eps A W[w;R].
\end{equation}
To justify differentiating the energy, truncate the symmetric double
potential integral and differentiate both moving domains. The terms with
$AR'$ cancel those with $j\cdot n$, leaving
$-\int_D\nabla K_\alpha w\cdot j$.
The potential and its first derivative have integrable majorants;
the potential is also integrable on the two-dimensional boundary since
$\alpha<1$. Truncations can thus be removed after time integration.
Finally
$\int_D\nabla w\cdot\nabla K_\alpha w
=\int wL_\alpha w-AW[w;R]$, proving
\eqref{eq:pressure-identity}.

\subsection{Uniform barriers and continuation}

Recall $\beta=1-\alpha$. The regularized boundary layer must respond to the singularity
$d^{-\alpha}$ in \eqref{eq:indicator}. For $\kappa>0$ define
\begin{equation}\label{eq:layer}
 \ell_\eps(d)=\int_0^d y_\eps(z)\,\dd z,
 \qquad y_\eps(d)=\frac1\eps\int_0^\infty
 e^{-\kappa h/\eps}(d+h)^{-\alpha}\,\dd h.
\end{equation}
The function is increasing and concave, and
\begin{align}
 \kappa\ell_\eps'-\eps\ell_\eps''&=d^{-\alpha},
 &\ell_\eps(d)&\leq\frac{d^\beta}{\kappa\beta},\label{eq:layer-eq}\\
 \ell_\eps'(0)&=\Gamma(1-\alpha)\kappa^{\alpha-1}\eps^{-\alpha}.
                                                        \label{eq:layer-slope}
\end{align}
These follow by differentiation and the Gamma integral. The decreasing
derivative gives
$|\ell_\eps(d+h)-\ell_\eps(d)|\leq\ell_\eps(|h|)$,
so the $C^\beta$ seminorm is uniform. The exponential approximate identity
in \eqref{eq:layer} gives uniform convergence on bounded intervals to
$d^\beta/(\kappa\beta)$, by the integrable bound in
\eqref{eq:layer-eq} and this common modulus.

\begin{lemma}\label{lem:uniform-family}
After deleting a finite prefix of the seed sequence, there is a common
$T>0$ on which \eqref{eq:viscous}--\eqref{eq:stefan} have the local
realizations above. Constants independent of $j$ give
\begin{gather}
 A/2\leq w_j\leq U\quad\hbox{in }D_j(t),\qquad
 R_0/2<R_j(t)<2R_0,\qquad |R_j'|\leq K,\label{eq:uniform-basic}\\
 \max\{c(t),A-K\ell_{\eps_j}(R_j-r)\}
 \leq w_j(t,r)\leq
 \min\{U,A+K\ell_{\eps_j}(R_j-r)\},\label{eq:barriers}\\
 |\eps_j\partial_nw_j|\leq K\eps_j^\beta.\label{eq:boundary-flux}
\end{gather}
Here $c(0)=3A/4$ and $A/2\leq c(t)\leq3A/4$.
Constants may depend on the fixed $\alpha$, profile and seed collar bounds,
but not on the interior seed gradients.
\end{lemma}

\begin{proof}
Initially stop a regular solution when its radius leaves $[R_0/2,2R_0]$
or $|\eps\partial_nw|$ reaches a small constant $J_0$.
Mass conservation, the cap, and \eqref{eq:boundary-positive} give
$v[w](R)\geq v_*>0$, uniformly before this stopping time.
The velocity is uniformly bounded and has the common spatial modulus
\eqref{eq:velocity-holder}. For a collar state $z\geq A/2$,
\begin{align*}
 az^{a-1}v[w](r)-R'
 \leq{}&-[1-a2^{1-a}]A^{a-1}v[w](R)\\
 &+a(A/2)^{a-1}|v[w](r)-v[w](R)|+J_0/A.
\end{align*}
The margin $1-a2^{1-a}=1-1/\sqrt2$ is positive. Choose the collar width
$\delta_0$ and then $J_0$ so that the last two errors each take at most
one quarter of this margin. For some $\kappa>0$ this gives
\begin{equation}\label{eq:collar-drift}
 az^{a-1}v[w](r)-R'\leq-2\kappa
 \quad (R-\delta_0<r<R).
\end{equation}

Consider the two functions in \eqref{eq:barriers}. The lower one is initially
below $g_j\geq A$. A sufficiently large $K$ makes the upper one dominate
the common outer $C^1$ collar, and outside that collar it is the cap $U$.
This uses the uniform positive lower bound for $\ell_\eps$ at each fixed
positive distance and its concavity, which bound $d/\ell_\eps(d)$ there.
Subtract $A\mathbf1_D$ from either zero-extended barrier. The remainder
is globally $C^\beta$ with norm at most a constant times $K$.
Since $\beta>\alpha$, \eqref{eq:indicator} and the H\"older estimate for
$L_\alpha$ give
\[
 |L_\alpha(\hbox{barrier}-A\mathbf1_D)|\leq K_2K,
 \qquad 0\leq L_\alpha\mathbf1_D\leq K_2d^{-\alpha}.
\]
At a smooth contact in the lower collar, \eqref{eq:collar-drift} and
\eqref{eq:layer-eq} bound the residual of
$\partial_t+a z^{a-1}v[w]\cdot\nabla+z^aL_\alpha-\eps\Delta$
by
\[
 -K d^{-\alpha}+K_2U^aA d^{-\alpha}+K_2U^aK.
\]
The radial curvature term is favorable. First enlarge $K$ to absorb the
$A$ term, then reduce $\delta_0$ so that $d^{-\alpha}$ absorbs the last
term, and finally enlarge $K$ to place both nonconstant collars inside the
reduced $\delta_0$. The fractional bound is linear in $K$, so the last
choice preserves the sign. The upper collar has reversed favorable
transport and layer signs and is a supersolution by the same choices.

The constant lower branch lies a fixed positive distance from the boundary:
at its join, $A-c\geq A/4$ and
$\ell_\eps(d)\leq d^\beta/(\kappa\beta)$. At a minimum of this barrier
only exterior vacuum makes a positive contribution to its fractional
Laplacian; it is at most $K_3c$ with a uniform $K_3$.
Choose $c'=-K_1c$ to dominate this term. Then $c\geq A/2$ on a common
positive time interval. The upper constant branch is a global maximum.
The lower join is convex and the upper join concave, so the incompatible
one-sided slopes preclude a forbidden first contact there.

Comparison uses the actual field $v[w]$ throughout. At a lower smooth
contact the zero-extended difference has global minimum zero and
$L_\alpha(w-w_-)\leq0$; at an upper contact the sign is reversed.
For fixed $\eps$, the difference equation has a zeroth-order coefficient
bounded below: the barrier gradient is bounded at that $\eps$, its density
is positive, and its fractional Laplacian is bounded below by $-K_2K$.
An exponential time factor therefore gives the non-strict minimum
principle. This proves comparison without assuming monotonicity of the
map $w\mapsto v[w]$.

Dividing the collar inequalities by $d$ at the boundary and using
\eqref{eq:layer-slope} gives \eqref{eq:boundary-flux}, which is strictly
below $J_0$ for sufficiently small viscosities. Thus the flux stopping
condition cannot occur first. Equation \eqref{eq:stefan} bounds $R'$;
shortening the common interval prevents either radius exit or $c=A/2$.

It remains to exclude an earlier loss of the fixed-$\eps$ realization.
Recall the fixed-ball variables $y=x/R(t)$,
$z(\tau,y)=w(t,R(t)y)$, $\rho(\tau)=R(t)$ and
$\dd\tau/\dd t=\eps/R(t)^2$.
The boundary derivative in \eqref{eq:forcing} is bounded at fixed $\eps$
by \eqref{eq:boundary-flux}. The other drift coefficients are bounded by
the density cap, positive floor and radius bounds. The fractional map of
$(z-A)\mathbf1_{B_1}$ is linear and bounded from $C^1$ to $L^\infty$.
Consequently
\[
 \|\mathcal N(z,\rho)\|_{p_0}\leq K_\eps(1+\|z\|_{C^1}).
\]
Equation \eqref{eq:heat-smoothing} gives a Volterra inequality with kernel
$1+r^{-\lambda}$, $\lambda<1$. Split each bounded rescaled interval into
finitely many pieces on which
$K_\eps\int_0^\delta(1+r^{-\lambda})\dd r<1/2$;
successive absorption bounds the $C^1$ norm on the whole interval.
The contraction can be restarted from positive $C^1$ data using these same
bounds and heat comparison. At positive elapsed times, the increment
estimate gives H\"older forcing. Applying $G$ to Duhamel after subtracting
the terminal forcing value is legitimate since
$\int_0^h r^{-1}r^\theta\dd r<\infty$ for $\theta>0$; the constant part
has generator $(\Sigma(h)-I)\mathcal N(z(t),\rho(t))$.
Closedness of $G$ gives the strict realization there, and interior heat
localization gives its classical derivatives. Local uniqueness matches
this restart to the earlier strict solution. Finite physical time is
bounded rescaled time at fixed $\eps$ because the radius stays between
positive bounds. Hence each fixed-$\eps$ solution continues to the common
physical time. No $C^1$ bound uniform in $\eps$ is asserted.
\end{proof}

\subsection{A compact limit and its global balances}

Taking $q=2-a$ in \eqref{eq:entropy-global} makes
$\Psi_{k^q}$ a positive multiple of $k$; taking $q=2$ controls the
viscous gradient. The boundary errors are $O(\eps_j^\beta)$ and the
outer weights are uniformly bounded. Thus
\begin{equation}\label{eq:compactness-bounds}
 \sup_j\int_0^T\|w_j\|_{H^{\alpha/2}(\R^3)}^2\,\dd t<\infty,
 \qquad \sup_j\eps_j\int_0^T\int_{D_j(t)}|\nabla w_j|^2<\infty.
\end{equation}
The pressure flux is bounded on a common enclosing ball. The zero-extended
viscous flux tends to zero in $L^2$, and \eqref{eq:extended-mass} gives
\begin{equation}\label{eq:time-bound}
 \partial_tw_j\ \hbox{bounded in }L^2_tH^{-1}_x,
 \qquad \|w_j(t)-g_j\|_{H^{-1}}\leq K(t+\sqrt{\eps_jt}).
\end{equation}
Here $H^{-1}$ is the dual of $H^1_0$ on a slightly larger fixed ball;
no whole-space $H^1$ norm across the outer jump is used.

The compactness lemma \cite[Lemma~7.5]{Stan2016} applies with
$X=H^{\alpha/2}$, $B=L^2$, $Y=H^{-1}$ and exponent two. Rellich gives
the first compact embedding, and Poincar\'e gives the second continuous
one. The derivative bound implies the required $L^1$-in-time bound on
finite intervals. A subsequence therefore satisfies
\begin{equation}\label{eq:limit-convergence}
 w_j\to u\quad\hbox{strongly in }L^2([0,T]\times\R^3),
 \qquad R_j\to R\quad\hbox{uniformly}.
\end{equation}
For smooth spatial tests, \eqref{eq:time-bound} gives equicontinuity.
A diagonal extraction on a countable dense family selects a density
$u(t)$ at every time, with $0\leq u(t)\leq U$. Approximate any bounded
spatial test in $L^1$ on the common support by a smooth test; the cap
bounds both approximation errors uniformly in time. Hence
\begin{equation}\label{eq:weak-time}
 \sup_{0\leq t\leq T}\left|\int(w_j(t)-u(t))\psi\right|\to0
 \quad(\psi\in L^\infty),\qquad
 u(0)=F,\quad \int u(t)=\int F.
\end{equation}
All such pairings are continuous. The uniformly convergent radii imply
$u(t)=0$ outside $B_{R(t)}$ at every time.

Truncate the near singularities of the potential and velocity. The retained
kernels on the common support form a compact family of spatial $L^1$
tests as the evaluation point varies on a compact set. A finite cover,
\eqref{eq:weak-time} and the near-field bounds give
\begin{equation}\label{eq:field-limit}
 K_\alpha w_j\to K_\alpha u,\qquad v[w_j]\to v[u]
 \quad\hbox{uniformly on compact space-time sets}.
\end{equation}
The Stefan law and \eqref{eq:boundary-flux} then yield
\begin{equation}\label{eq:radius-limit}
 R_j\to R\quad\hbox{in }C^1([0,T]),\qquad
 R'=A^{a-1}v[u](R),\qquad R(0)=R_0.
\end{equation}
The layer limit gives $u\geq c(t)\geq A/2$ inside its ball and
\begin{equation}\label{eq:trace-modulus}
 |u(t,r)-A|\leq\frac{K}{\kappa\beta}(R(t)-r)^\beta
 \quad\hbox{near }R(t).
\end{equation}
These inequalities pass by bounded-test pairings, so hold at every time
in the essential spatial sense. In particular the inner trace is $A$.
Strong density convergence and \eqref{eq:field-limit} pass the nonlinear
flux in local $L^1$; the viscous flux vanishes. The same $u$ solves the
whole-space equation with initial $F$ and locally integrable flux.

The energies in \eqref{eq:pressure-identity} converge at every terminal
time by \eqref{eq:weak-time}, \eqref{eq:field-limit} and
$\mathcal E_\alpha(w)=\tfrac12\int wK_\alpha w$.
The pressure dissipation converges in $L^1$ on time intervals by bounded
fields and strong density convergence. Its positive error
$\eps_j A\int_0^tW[w_j;R_j]\dd r$ vanishes, while the other viscous term
is nonnegative. Therefore
\begin{equation}\label{eq:limit-pressure}
 \mathcal E_\alpha(u(t))+
 \int_0^t\int u^a|\nabla K_\alpha u|^2
 \leq\mathcal E_\alpha(F),\qquad 0\leq t\leq T.
\end{equation}

To retain the viscosity in the power balances, extract once the joint
nonnegative finite measures
\begin{equation}\label{eq:nu}
 \nu_j=\eps_j|\nabla w_j|^2\mathbf1_{D_j(t)}\,\dd x\dd t\,
             \delta_{w_j(t,x)}(\dd k)\ \rightharpoonup\ \nu.
\end{equation}
Their mass is bounded by \eqref{eq:compactness-bounds}; they have common
compact space-time support and state support $[A/2,U]$.
Weak compactness follows, for example, by diagonal extraction on a
countable dense family of continuous functions on that compact set and
extension of the resulting positive bounded functional. This extraction
is independent of the entropy exponent.

At good terminal times, \eqref{eq:entropy-global} passes to
\begin{align}
 &\|u(t)\|_q^q+\int_0^t\mathcal B_\alpha(\Psi_{k^q},u(r))\,\dd r\nonumber\\
 &\quad+q(q-1)\int_{(0,t)\times\R^3\times[A/2,U]}k^{q-2}\,\dd\nu\nonumber\\
 &\leq\|F\|_q^q+q(q-1)\int_0^tW[u(r);R(r)]\,\dd r
             \int_0^A k^{q-2}g_A(k)\,\dd k.\label{eq:limit-power-cost}
\end{align}
A good time has strong spatial convergence along the subsequence and no
time atom of $\nu$. Fatou gives the positive fractional term, measure
convergence gives the state moment, and \eqref{eq:field-limit} passes the
surface weight. The boundary error vanishes. Take first a countable dense
set of $q>1$ and intersect the good-time sets. Extension to real $q$ uses
bounded endpoint powers, continuity of $k^{q-2}$ on $[A/2,U]$ and Fatou
for the positive pair integrands. Near zero, the outer density integrals
are dominated by a multiple of
$k^{a+q_{\min}-2}+k^{q_{\min}-1}$ in a neighborhood of each $q>1$.

In this passage initially include time zero in the measure moment. The
quadratic case gives, at good times,
\[
 \|u(t)\|_2^2\leq\|F\|_2^2+Kt,\qquad
 2\nu([0,t]\times\R^3\times[A/2,U])
 \leq\|F\|_2^2-\|u(t)\|_2^2+Kt.
\]
Weak-time lower semicontinuity extends the first bound to every time.
The initial weak trace and cap imply weak $L^2$ convergence to $F$ on
the common support. The norm bound makes it strong as $t\downarrow0$.
The second bound now excludes an initial atom of $\nu$, justifying
$(0,t)$ in \eqref{eq:limit-power-cost}.

There is a fixed annulus about $S$ on which the velocity remains positive
for a common interval beginning at zero, uniformly in the retained
approximants. Indeed \eqref{eq:time-bound} controls smooth retained
velocity kernels; their omitted near kernels are uniformly small. A finite
test cover gives a common time modulus near $S$, and the initial fields
converge uniformly there to $v[F]$. Its value at $S$ is positive and its
radial component is continuous. Shorten $T$ and discard a finite prefix
so that an annulus $J_+$, separated from zero and all outer boundaries,
satisfies
\begin{equation}\label{eq:positive-annulus}
 v[w_j](t,r),\ v[u](t,r)\geq v_*>0,
 \qquad (t,r)\in[0,T]\times J_+.
\end{equation}
Finite deletion changes neither limit. A restriction time with no terminal
atom of $\nu$ can be chosen. We use this same restricted density, radius,
tail sequence and measure in the rest of the proof.

\section{Interior dissipation}\label{sec:interaction}

The global balance \eqref{eq:limit-power-cost} contains a positive boundary
entropy contribution. To compare it with dissipation at the internal jump, we must separate
the viscosity measure from the nonlocal production. The latter will have
an $L^2$ density near $S$ and therefore no contribution at the scale of the
initial shock.

For a normalized smooth entropy, $\eta(0)=\eta'(0)=0$. Recall
$\Psi_\eta'=\eta''f$, $F_\eta'=\Psi_\eta$, with zero constants at zero,
and $Q_\eta=\eta'f-\Psi_\eta$. Define
\begin{equation}\label{eq:defect}
 \mathcal D_\eta[w](x)=c_{3,\alpha}\int
 \frac{F_\eta(w(y))-F_\eta(w(x))
 -\Psi_\eta(w(x))(w(y)-w(x))}{|x-y|^{3+\alpha}}\,\dd y.
\end{equation}
It equals $\Psi_\eta(w)L_\alpha w-L_\alpha F_\eta(w)$ wherever the
terms are realized by their integrals. Since $F_\eta''=\eta''f$, it is
nonnegative for convex $\eta$. On an interior annulus the regularized
equation gives the two exact radial balances
\begin{align}
 \partial_t(r^2w)+\partial_r(r^2v f(w)-\eps r^2w_r)&=0,
                                                        \label{eq:radial-mass}\\
 \partial_t(r^2\eta(w))+
 \partial_r(r^2v Q_\eta(w)-\eps r^2\eta'(w)w_r)
 &=-r^2\bigl(\eps\eta''(w)w_r^2+\mathcal D_\eta[w]\bigr)
       -rL_{1,\alpha}[rF_\eta(w)].\label{eq:radial-entropy}
\end{align}
Here the profile $rF_\eta(w)$ is extended oddly to the real line.
The identity behind the last term is
\begin{equation}\label{eq:radial-operator}
 rL_{3,\alpha}h=L_{1,\alpha}[rh]
\end{equation}
for bounded compactly supported radial $h$. For smooth $h$, the
three-dimensional radial Fourier transform is a constant times
$\rho^{-1}\int_0^\infty rh(r)\sin(r\rho)\dd r$; the odd
one-dimensional transform has the same sine transform. Multiplication
by $|\rho|^\alpha$ proves \eqref{eq:radial-operator}, with the same
normalization on both sides. Radial mollification passes it to bounded
compact $h$: the odd profiles converge in $L^1(\R)$, since their integral
near zero is uniformly $O(\delta^2)$ and away from zero radial weighted
$L^1$ convergence suffices. Transfer the fractional operator to the smooth
test to pass to distributions. The boundedness assumption here includes
the actual profiles $F_\eta(w)$.

For $\eta(k)=k^2/2$, the transforms are
$\Psi_\eta(k)=k^{a+1}/(a+1)$ and
$F_\eta(k)=k^{a+2}/((a+1)(a+2))$.
The total masses of $\eps|\nabla w|^2$ and
$\mathcal D_\eta[w]$ are uniformly bounded. Indeed the former is in
\eqref{eq:compactness-bounds}, while integration of the latter over space
is $\mathcal B_\alpha(\Psi_\eta,w)$: symmetric truncation gives
$\int L_\alpha F_\eta(w)=0$. Equation \eqref{eq:entropy-global}
then bounds it. In radial coordinates these measures have the weight
$r^2$ and the full-space integrals are $4\pi$ times the radial ones.

\subsection{A weighted interaction estimate}

We use the interaction calculation of \cite[Section~2]{GolsePerthame2013},
with the weighted and fractional terms retained explicitly. If compactly
supported smooth fields satisfy
$A_t+B_r=C$ and $D_t+E_r=H$, set
\[
 I(t)=\int A(t,r)\int_r^\infty D(t,z)\,\dd z\dd r,
 \qquad P(t,r)=\int_{-\infty}^r A(t,z)\,\dd z.
\]
Differentiation and integration by parts give
\begin{equation}\label{eq:interaction-identity}
 \int_0^T\int(AE-DB)=I(T)-I(0)
 -\int_0^T\int C(t,r)\int_r^\infty D(t,z)\,\dd z\dd r\dd t
 -\int_0^T\int HP.
\end{equation}
The identity also holds for bounded measurable mass and entropy densities
and fluxes, bounded $C$, and a source
$H=\mu+\sum_l b_lL_{1,\alpha}g_l$, where $\mu$ is a finite signed
measure, $b_l$ are compactly supported Lipschitz functions, and $g_l$ are
bounded on the full line. The necessary interpretation is
\[
 \int HP=\int P\,\dd\mu+
             \sum_l\int g_lL_{1,\alpha}(b_lP).
\]
To verify it, the mass equation gives
$P_r=A$ and $P_t=-B+\int_{-\infty}^rC$, so $P$ has a locally Lipschitz
representative in both variables. Multiply by a spatial cutoff equal to
one on the entropy support and mollify. The primitives converge uniformly
on compact sets, their derivatives converge locally in $L^1$, and their
$W^{1,1}$ norms stay bounded. The fractional pairings converge by
\begin{equation}\label{eq:fractional-L1}
 \|L_{1,\alpha}\psi\|_1
 \leq K_\alpha\|\psi\|_1^{1-\alpha}\|\psi'\|_1^\alpha,
 \qquad \psi\in W^{1,1}(\R).
\end{equation}
This estimate follows by integrating
$\|\psi(\cdot+h)-\psi\|_1\leq
\min(2\|\psi\|_1,|h|\|\psi'\|_1)$ against $|h|^{-1-\alpha}$.
Uniform convergence handles the measure pairing. Thus
\eqref{eq:interaction-identity} holds at Lebesgue trace endpoints with no
source-measure atom, or with smooth time cutoffs. A prescribed bounded
initial mass and entropy trace is included if there is no initial atom;
approach zero through good times and retain the initial terms. Source
integrals between endpoints use the open time interval.

Write $\Delta_hz(r)=z(r+h)-z(r)$ and take a fixed smooth spatial cutoff
$\chi$. For a bounded field $z_0$ and a bounded Lipschitz weight $\omega$,
\[
 P_h(r)=\int_{-\infty}^r\chi(z)\Delta_h(\omega z_0)(z)\,\dd z
 \quad\hbox{satisfies}\quad
 \|P_h\|_\infty\leq K|h|,\qquad \|(P_h)'\|_\infty\leq K
\]
for bounded $z_0$. Move the translation to the integration endpoint and
the cutoff to obtain the first bound. For a uniformly Lipschitz coefficient
$b$ near the cutoff support, the compact test
\begin{equation}\label{eq:correct-test}
 \psi_h=b\,\Delta_{-h}(\chi P_h)
\end{equation}
satisfies $\|\psi_h\|_1\leq K|h|$ and
$\|\psi_h'\|_1\leq K$. Equation \eqref{eq:fractional-L1} therefore
bounds its pairing with $L_{1,\alpha}g$ by
$K\|g\|_\infty|h|^{1-\alpha}$. No $L^1$ bound on $g$ is needed.

\begin{lemma}\label{lem:increments}
Let bounded states $z_\eps\in[A/2,U]$ obey, on a larger time-interval
cylinder with spatial interval $J_+$,
\begin{align*}
 \partial_t(\omega z_\eps)+
 \partial_r(\omega vf(z_\eps)-\eps\omega\partial_rz_\eps)&=0,\\
 \partial_t(\omega\eta(z_\eps))+
 \partial_r(\omega vQ_\eta(z_\eps)-\eps\omega\eta'(z_\eps)\partial_rz_\eps)
 &=-\mu-bL_{1,\alpha}g,
 \qquad\eta(k)=k^2/2.
\end{align*}
Assume $\omega$ is time independent and uniformly positive, $\omega,b$
have bounded Lipschitz norms, $v\geq v_*>0$ is bounded and spatially
$C^\beta$, $\mu\geq0$ has bounded mass, and $g$ is uniformly bounded
on the full line. For $\eps>0$ assume local boundedness of
$\partial_rz_\eps$ for each member and
$\eps\int|\partial_rz_\eps|^2\leq K$; for $\eps=0$ no derivative is
required. At good endpoints as above, or with a fixed smooth time cutoff,
every smaller cylinder $Q$ satisfies
\begin{equation}\label{eq:L4-increment}
 \int_Q|\Delta_hz_\eps|^4\leq K_Q(|h|^\beta+\sqrt\eps).
\end{equation}
The constants depend on the listed bounds, intervals and cutoffs. The
support of $g$ may grow.
\end{lemma}

\begin{proof}
For $z_0,z_1\in[A/2,U]$, $Q_\eta'=\eta'f'$ gives
\begin{align}
 &(z_1-z_0)(Q_\eta(z_1)-Q_\eta(z_0))
 -(\eta(z_1)-\eta(z_0))(f(z_1)-f(z_0))\nonumber\\
 &\quad=\frac12\int_{z_0}^{z_1}\int_{z_0}^{z_1}
 (r-q)(f'(r)-f'(q))\,\dd r\dd q
 \leq-\frac{a(1-a)U^{a-2}}{12}|z_1-z_0|^4.\label{eq:determinant}
\end{align}
The formula and inequality hold in either ordering of the states, since
$f''(k)\leq-a(1-a)U^{a-2}$ on this interval.

Put $J_m=\omega vf(z_\eps)-\eps\omega\partial_rz_\eps$ and
$J_e=\omega vQ_\eta(z_\eps)-\eps\omega\eta'(z_\eps)\partial_rz_\eps$.
In \eqref{eq:interaction-identity} use
\[
 A=\chi\Delta_h(\omega z_\eps),\ B=\chi\Delta_hJ_m,
 \ C=\chi'\Delta_hJ_m,\quad
 D=\chi\Delta_h(\omega\eta(z_\eps)),\ E=\chi\Delta_hJ_e,
\]
and $H=\chi'\Delta_hJ_e-\chi\Delta_h\mu
-\chi\Delta_h(bL_{1,\alpha}g)$. The weak extension applies: at each
fixed viscous member the fluxes are locally bounded; at zero viscosity
bounded states and velocity suffice. The mass primitive has the bounded
space and time derivatives described above, so its measure pairing is
well defined. Both density primitives are $O(|h|)$ uniformly. Endpoints,
cutoff sources and the translated measure therefore cost $O(|h|)$,
apart from viscous fluxes, whose $L^1$ norms are $O(\sqrt\eps)$ by
Cauchy--Schwarz.

Move the fractional-source translation onto the primitive. The resulting
test is exactly $b\Delta_{-h}(\chi P_h)$, with $b$ outside the difference.
Equations \eqref{eq:fractional-L1} and \eqref{eq:correct-test} bound this
term by $K|h|^{1-\alpha}$ using only $\|g\|_\infty$.
Freezing $\omega,v$ at the unshifted point in $AE-DB$ leaves
$\chi^2\omega^2v$ times \eqref{eq:determinant}. Weight errors cost
$O(|h|)$ and velocity errors $O(|h|^\beta)$; the remaining viscous
products have integral $O(\sqrt\eps)$ because all other factors are
bounded. The positive lower bounds for $\omega$ and $v$ now give
\eqref{eq:L4-increment}. The same proof at $\eps=0$ has no viscosity
error, and the cutoff version retains the corresponding initial terms.
\end{proof}

\subsection{Separation of the nonlocal defect}

Apply Lemma~\ref{lem:increments} to \eqref{eq:radial-mass}--\eqref{eq:radial-entropy}
with quadratic entropy. The actual coefficients are
\[
 \omega=r^2,\qquad b=r,\qquad
 g=rF_\eta(w_j)\ \hbox{extended oddly},\qquad
 \mu=r^2\bigl(\eps_j|\partial_rw_j|^2+\mathcal D_\eta[w_j]\bigr)\,\dd r\dd t.
\]
The common annulus \eqref{eq:positive-annulus}, the weighted mass bounds
above and \eqref{eq:velocity-holder} verify all its hypotheses.

\begin{lemma}\label{lem:defect-convergence}
For each fixed normalized $C^2$ entropy $\eta$ with bounded second
derivative on $[0,U]$, convex or not,
\[
 \mathcal D_\eta[w_j]\longrightarrow\mathcal D_\eta[u]
 \quad\hbox{strongly in }L^2
\]
on each smaller time-annulus cylinder, including the initial endpoint.
In particular, for every fixed $K_1$,
\begin{equation}\label{eq:tube-defect}
 \frac1\lambda\int_0^\lambda
 \int_{|r-S|<K_1\lambda}|\mathcal D_\eta[u](t,r)|\,\dd r\dd t
 \longrightarrow0.
\end{equation}
The same statement holds with the full radial volume weight.
\end{lemma}

\begin{proof}
For $r,\rho$ in an annulus away from zero, angular integration of the
three-dimensional kernel is a positive constant times
\[
 \frac\rho r\left(|r-\rho|^{-1-\alpha}-(r+\rho)^{-1-\alpha}\right).
\]
Near its diagonal it is bounded by $K|r-\rho|^{-1-\alpha}$.
Taylor's integral remainder in \eqref{eq:defect} is bounded in absolute
value by $K_\eta|w_j(\rho)-w_j(r)|^2$, since
$F_\eta''=\eta''f$ is bounded on $[0,U]$.
Minkowski consequently bounds the $L^2$ norm of the contribution from
$|r-\rho|<\delta$ by
\begin{equation}\label{eq:diagonal-bound}
 K_\eta\int_0^\delta h^{-1-\alpha}\|\Delta_hw_j\|_4^2\,\dd h.
\end{equation}
The norms here are on a slightly larger time-annulus cylinder.
Bounded states and the viscous gradient estimate give
\[
 \|\Delta_hw_j\|_4^2\leq K\|\Delta_hw_j\|_2
 \leq K|h|/\sqrt{\eps_j}.
\]
Lemma~\ref{lem:increments} also gives
$\|\Delta_hw_j\|_4^2\leq K(h^{(1-\alpha)/2}+\eps_j^{1/4})$.
Split \eqref{eq:diagonal-bound} at $h_j=\eps_j^{2/3}$.
The first estimate below $h_j$ gives
$O(\eps_j^{1/6-2\alpha/3})$; the second above $h_j$ gives
\[
 K_\eta\delta^{(1-3\alpha)/2}
       +O(\eps_j^{1/4-2\alpha/3}).
\]
Both error exponents are positive because $\alpha<1/4$.
Thus the iterated near-diagonal bound tends to zero with $\delta$,
also on intervals abutting zero.

Away from the diagonal the kernel has bounded integrable operator norms.
The remainder depends Lipschitz continuously on its two states in
$[0,U]$. Strong density convergence therefore gives strong $L^2$
convergence of this truncated defect, including its separated contribution
from the origin and exterior vacuum. Fatou gives the same near-diagonal
bound for the limit. Let $j\to\infty$ and then $\delta\downarrow0$ to
obtain strong $L^2$ convergence and local $L^2$ integrability of the limit.

The tubes in \eqref{eq:tube-defect} lie in one such fixed cylinder about
$S$ and have area $O(\lambda^2)$. Cauchy--Schwarz bounds their integral
divided by $\lambda$ by a constant times the $L^2$ norm on the tube,
which tends to zero by absolute continuity. The volume weight $4\pi r^2$
is bounded there, proving the full-space version.
\end{proof}

We may now pass the local entropy equation without a chain rule for the
weak limit. The viscous entropy flux tends to zero in $L^2$, since
$\eta'$ is bounded on $[A/2,U]$ and
$\eps_j\int|\nabla w_j|^2$ is bounded. The moment of $\eta''$ in
\eqref{eq:nu} converges. Lemma~\ref{lem:defect-convergence} passes the
nonlocal defect; the linear fractional term passes by transferring the
operator to the test. Density and other fluxes pass strongly, and the
initial entropies pass from $g_j\to F$. Therefore, on each smaller
time-annulus cylinder contained in $[0,T]\times J_+$,
\begin{equation}\label{eq:limit-local-entropy}
 \partial_t\eta(u)+\nabla\cdot(v[u]Q_\eta(u))+L_\alpha F_\eta(u)
 =-\int\eta''(k)\nu(\dd t,\dd x,\dd k)
   -\mathcal D_\eta[u]\,\dd t\dd x,
\end{equation}
with initial term $\eta(F)$. The notation on the right means the
state-variable moment of the single measure $\nu$ already extracted.
It holds for each fixed normalized $C^2$ entropy with bounded second
derivative, by the same argument, including nonconvex ones.

Rotational invariance of the approximants implies rotational invariance
of $\nu$. Define its radial measure by
\begin{equation}\label{eq:radial-nu}
 \bar\nu=\frac1{4\pi}(t,x,k\mapsto t,|x|,k)_\#\nu.
\end{equation}
The radial version of \eqref{eq:limit-local-entropy} has precisely
$-\int\eta''(k)\bar\nu$ and
$-r^2\mathcal D_\eta[u]\dd r\dd t-rL_{1,\alpha}[rF_\eta(u)]\dd r\dd t$
as its sources. Thus the radial measure is determined by the full-space
one; it is not another independently selected limit. The strong initial
trace and absence of an initial atom established above justify the same
initial terms in this radial equation.

\section{The initial shock and all power inequalities}\label{sec:shock}

We examine the internal jump at $r=S$, where the velocity is positive.
For $\lambda\downarrow0$ put
\begin{equation}\label{eq:tangent}
 U_\lambda(\tau,\xi)=u(\lambda\tau,S+\lambda\xi),\quad
 r_\lambda=S+\lambda\xi,\quad
 v_\lambda=v[u](\lambda\tau,r_\lambda),\quad v_0=v[F](S)>0.
\end{equation}
On every fixed half-cylinder, $r_\lambda^2\to S^2$ and
$v_\lambda\to v_0$ uniformly, and $A/2\leq U_\lambda\leq U$ for
small $\lambda$. The initial data tend in every finite local $L^q$ to
the step $d_*(\xi)=1$ for $\xi<0$ and $d_*(\xi)=U$ for $\xi>0$.

\subsection{Compactness and identification of the tangent}

The scaled mass flux is $r_\lambda^2v_\lambda f(U_\lambda)$ and the
scaled entropy flux is $r_\lambda^2v_\lambda Q_\eta(U_\lambda)$.
For the linear fractional term, let
\[
 \widetilde g(t,r)=rF_\eta(u(t,|r|))\quad(r\in\R),\qquad
 g_\lambda(\tau,\xi)=\widetilde g(\lambda\tau,S+\lambda\xi).
\]
It is the pullback of the odd physical radial profile. Its supremum norm
is bounded uniformly, although its support in $\xi$ grows. The scaled term is
\begin{equation}\label{eq:scaled-source}
 \lambda^{1-\alpha}r_\lambda L_{1,\alpha}g_\lambda.
\end{equation}
Transferring $L_{1,\alpha}$ to a compact smooth test bounds its pairing by
$\lambda^{1-\alpha}\|g_\lambda\|_\infty
\|L_{1,\alpha}(r_\lambda\varphi)\|_1\to0$.
The last norm is uniformly bounded by \eqref{eq:fractional-L1}.

For a convex entropy the other scaled sources are positive measures.
Testing the entropy balance with a nonnegative cutoff equal to one on a
fixed rectangle bounds their mass there: the entropy, flux and initial
terms are bounded, and \eqref{eq:scaled-source} has a bounded pairing.
In particular the quadratic measure has a uniform local mass bound.
The exact coefficients for Lemma~\ref{lem:increments} at zero viscosity are
\[
 \omega_\lambda=r_\lambda^2,\qquad
 b_\lambda=\lambda^{1-\alpha}r_\lambda,\qquad
 g=g_\lambda,\qquad\eps=0.
\]
The weights are uniformly positive and Lipschitz on each fixed rectangle.
The spatial H\"older seminorm of $v_\lambda$ gains $\lambda^\beta$ from
\eqref{eq:velocity-holder}; no time derivative of $v_\lambda$ is required.
Thus, on a smaller fixed rectangle,
\begin{equation}\label{eq:scaled-increments}
 \int|U_\lambda(\tau,\xi+h)-U_\lambda(\tau,\xi)|^4\,\dd\xi\dd\tau
 \leq K|h|^{1-\alpha}.
\end{equation}
This is a fourth-power bound; the corresponding $L^4$ norm has exponent
$(1-\alpha)/4$. It uses the uniform supremum of $g_\lambda$, without
a bound on its growing spatial $L^1$ norm.

Cauchy--Schwarz turns \eqref{eq:scaled-increments} into a squared
$L^2$ increment bound $K|h|^{(1-\alpha)/2}$. After spatial localization,
the translation formula for Sobolev norms gives a uniform
$L^2_\tau H^\sigma_\xi$ bound for each
$0<\sigma<(1-\alpha)/4$. The weighted mass equation bounds
$\partial_\tau(r_\lambda^2U_\lambda)$ in $L^1_\tau H^{-1}_\xi$.
Use \cite[Lemma~7.5]{Stan2016} on bounded intervals with
$X=H^\sigma$, $B=L^2$, $Y=H^{-1}$ and exponent two, then divide by the
smooth positive weight. A diagonal extraction gives, from every scale
sequence, a tangent $Z$ with strong local $L^2$ convergence.
If compactness is first applied away from $\tau=0$, the cap discards an
initial strip with squared $L^2$ mass $O(\delta)$ uniformly, so the
convergence extends to compact half-cylinders. Bounded states upgrade it
to every finite local $L^q$.

Passing the scaled mass and convex entropy balances yields
\begin{equation}\label{eq:scalar-limit}
 Z_\tau+v_0\partial_\xi f(Z)=0,\qquad
 \partial_\tau\eta(Z)+v_0\partial_\xi Q_\eta(Z)\leq0,
\end{equation}
with the actual step initial mass and entropy terms. Subtracting affine
functions extends the normalized inequalities to all smooth convex
entropies. The mass equation selects compact-test continuous pairings.
The localized quadratic entropy inequality gives
\[
 \limsup_{\tau\downarrow0}\int\chi Z(\tau)^2
 \leq\int\chi d_*^2
\]
for nonnegative compact $\chi$, first at good times and then by weak
lower semicontinuity for the representative. Its weak initial trace and
bounded states give $\int\chi Z(\tau)d_*\to\int\chi d_*^2$, by smooth
approximation of $\chi d_*$. Expanding the square proves the strong
local $L^2$, hence local $L^1$, initial trace. This supplies the scalar
initial trace without assuming a strong trace at the unscaled inner jump.

The bounded scalar uniqueness theorem
\cite[Section~6.1, Theorem~4]{Bressan2009} now identifies the tangent.
To match its smooth-flux hypothesis, extend $v_0k^a$ smoothly to $\R$
outside a neighborhood of $[A/2,U]$, using a cutoff supported in the
positive half-line. All compared states stay in that interval, so the
equation and entropy flux derivatives are unchanged. The step itself
need not be integrable; the theorem treats bounded data, and the compared
initial difference is zero. The solution is
\begin{equation}\label{eq:shock}
 Z_*(\tau,\xi)=\begin{cases}1,&\xi<\sigma_0\tau,\\
 U,&\xi>\sigma_0\tau,\end{cases}
 \qquad \sigma_0=v_0\frac{U^a-1}{U-1}=\frac{v_0}{11}.
\end{equation}
It satisfies Rankine--Hugoniot. Recall the chord deficit
$H_U(k)=f(k)-f(1)-[(f(U)-f(1))/(U-1)](k-1)$.
For every convex $C^2$ entropy, integration by parts in $k$, using
$H_U(1)=H_U(U)=0$, gives its exact
entropy distribution
\begin{equation}\label{eq:shock-production}
 \partial_\tau\eta(Z_*)+v_0\partial_\xi Q_\eta(Z_*)
 =-v_0\left(\int_1^U\eta''(k)H_U(k)\,\dd k\right)
             \delta_{\xi=\sigma_0\tau}\,\dd\tau.
\end{equation}
Its sign is nonpositive. Uniqueness identifies every subsequential tangent
with \eqref{eq:shock}. Since every scale sequence is relatively compact,
the entire family $U_\lambda$ converges to $Z_*$ in every finite local
$L^q$.

\subsection{Dissipation at higher densities}

Fix any $A<B<1$. Choose $0\leq\zeta\leq1$ smooth, zero for $k\leq B$
and one on $[1,U]$, and define a normalized convex entropy by
$\eta''(k)=\zeta(k)/k$, with zero extension near zero. It has bounded
second derivative on $[0,U]$, so \eqref{eq:limit-local-entropy} applies.
Its shock coefficient in \eqref{eq:shock-production} is $-v_0J$, where
$J=\int_1^U H_U(k)/k\,\dd k$.

Under $t=\lambda\tau$, $r=S+\lambda\xi$, the radial viscosity measure is
\[
 \bar\nu_\lambda=\lambda^{-1}
 (t,r,k\mapsto t/\lambda,(r-S)/\lambda,k)_\#\bar\nu.
\]
The factor $\lambda^{-1}$ follows from the first derivative scaling
and the two-dimensional time-radius Jacobian. In the exact entropy
balance, the linear fractional term vanishes by
\eqref{eq:scaled-source}, while \eqref{eq:tube-defect} makes the separated
nonlocal defect $o(\lambda)$ on each physical tube of width and length
$O(\lambda)$. Full tangent convergence therefore gives, weakly against
compact tests with $\tau>0$,
\begin{equation}\label{eq:viscosity-shock}
 \int\frac{\zeta(k)}k\,\bar\nu_\lambda(\dd\tau,\dd\xi,\dd k)
 \ \rightharpoonup\ S^2v_0J\,
              \delta_{\xi=\sigma_0\tau}\,\dd\tau.
\end{equation}
The local positive measure bounds above justify this weak interpretation.
In particular the nonlocal defect has not been counted as viscosity
production.

Test \eqref{eq:viscosity-shock} by a nonnegative cutoff at most one,
supported in $0<\tau<1$ and a fixed spatial interval containing the
shock segment, and equal to one on the shock for
$\delta\leq\tau\leq1-\delta$. Since
$\zeta(k)/k\leq\mathbf1_{[B,U]}(k)/k$, positivity and
\eqref{eq:radial-nu} imply, on sending $\delta\downarrow0$,
\begin{equation}\label{eq:high-density}
 \liminf_{t\downarrow0}\frac1t
 \int_{(0,t)\times\R^3\times[B,U]}\frac1k\,\dd\nu
 \geq4\pi S^2v_0J=W[F;S]J.
\end{equation}
This is a full liminf, since the tangent convergence used no favorable
scale sequence.

The outer weight $W[u(t);R(t)]$ is positive and continuous at zero, with
value $W[F;R_0]$. Recall that $G_A=\int_0^A g_A(k)/k\,\dd k$ and
the fixed strict margin is $W[F;R_0]G_A<W[F;S]J$. This margin and
\eqref{eq:high-density} therefore supply one $0<T_1<T$ such that
\begin{equation}\label{eq:payment}
 \int_{(0,t)\times\R^3\times[B,U]}\frac1k\,\dd\nu
 \geq G_A\int_0^tW[u(r);R(r)]\,\dd r,
 \qquad0<t\leq T_1.
\end{equation}
For example choose an error smaller than one third of the positive gap in
\eqref{eq:strict-margin}, apply the definition of liminf on one interval
and continuity of the averaged outer weight on another, and intersect
them. This time is chosen independently of every later power exponent.

\subsection{The exact power coefficient and every terminal time}

For $b>0$ and $r,z\geq0$, Cauchy--Schwarz on the interval between the
two states gives
\begin{equation}\label{eq:scalar-power}
 (r^b-z^b)(r-z)\geq\frac{4b}{(b+1)^2}
       (r^{(b+1)/2}-z^{(b+1)/2})^2.
\end{equation}
Indeed apply it to $\int_z^r k^{(b-1)/2}\dd k$, and then use a positive
lower endpoint and monotone convergence if one state is zero.
Set $b=a+q-1$, so
$\Psi_{k^q}(k)=q(q-1)k^b/b$ by \eqref{eq:psi-power}.
Combining this formula, \eqref{eq:scalar-power} and the positive pair form gives
\begin{equation}\label{eq:power-fractional}
 \mathcal B_\alpha(\Psi_{k^q},w)
 \geq\frac{4q(q-1)}{(a+q)^2}
 \|(-\Delta)^{\alpha/4}w^{(a+q)/2}\|_2^2.
\end{equation}
This holds as a comparison of nonnegative extended integrals. For
$q=2-a$ the scalar inequality is equality; no strict fractional remainder
is available at that exponent.

Take, for example, $B=(A+1)/2$. On $k\geq B$,
$k^{q-2}\geq B^{q-1}/k$. The viscosity moment in
\eqref{eq:limit-power-cost} is therefore at least $B^{q-1}$ times the
left side of \eqref{eq:payment}. The positive boundary contribution in that balance is at most
$A^{q-1}G_A\int_0^tW[u;R]$ by \eqref{eq:density-cost}.
Both terms carry the same factor $q(q-1)$ in \eqref{eq:limit-power-cost}.
Since $B>A$, the same measure and the same $T_1$ make the viscosity
contribution dominate the boundary contribution for every $q>1$.
At good times the result, followed by
\eqref{eq:power-fractional}, is
\begin{equation}\label{eq:short-power}
 \|u(t)\|_q^q+\frac{4q(q-1)}{(a+q)^2}
 \int_0^t\|(-\Delta)^{\alpha/4}u^{(a+q)/2}\|_2^2\,\dd r
 \leq\|F\|_q^q.
\end{equation}
For any terminal $t\leq T_1$, approach it by increasing good times.
Bounded-test weak $L^1$ continuity gives lower semicontinuity of the
convex power norm, and the nonnegative time integral increases to its
full value. This extends \eqref{eq:short-power} to every terminal time.
The dense-exponent passage described after \eqref{eq:limit-power-cost}
gives every real $q>1$ without changing the measure or time interval.
Together with \eqref{eq:limit-pressure}, the distributional equation,
mass, cap and \eqref{eq:weak-time}, this proves all required energy weak
solution properties on $[0,T_1]$. Its support is contained in the finite
$C^1$ radius $R(t)$ there.

\section{Global continuation}\label{sec:continuation}

The construction so far is on a positive finite interval. The question
requires a global solution, with every bounded spatial pairing continuous
and both kinds of inequality at every terminal time. We give the
representative argument needed to use an existing global existence theorem.

\begin{proposition}\label{prop:global-source}
For the fixed $0<\alpha<1/4$ and every nonzero nonnegative compactly
supported $d\in L^1(\R^3)\cap L^\infty(\R^3)$, there is a global
solution in the energy weak class of Definition~\ref{def:class}, with
$m=1+a=3/2$ and $s=1-\alpha/2$.
\end{proposition}

\begin{proof}
Apply \cite[Definition~2.1 and Theorem~2.2]{Stan2016} with
$N=3$, $m=3/2$ and pressure parameter $s=1-\alpha/2$.
It supplies a nonnegative global distributional solution $z$, its initial
compact-test formulation, locally integrable pressure gradient and flux,
boundedness and integrability on finite cylinders, conserved mass
$M=\int d$, cap $U_d=\|d\|_\infty$, and the pressure and every-$q$
power inequalities with coefficient $4q(q-1)/(a+q)^2$.
The source differentiating order is $1-s=\alpha/2$, so its seminorm is
exactly $\|(-\Delta)^{\alpha/4}z^{(a+q)/2}\|_2$.
In this dimension $3>2s$, the scalar inverse pressure is well defined
with \eqref{eq:riesz}; its composed gradient equals $\nabla K_\alpha z$.
The initial energy is finite for bounded compact $d$, as proved above.
We interpret time-slice estimates initially on their good-time set and
prove the stronger representative assertions next.

Split the kernel $k(x)=c_\alpha x/|x|^{3+\alpha}$ into
$k_{\rm near}=k\mathbf1_{|x|<1}$ and
$k_{\rm far}=k\mathbf1_{|x|\geq1}$. The radial integrals show
$k_{\rm near}\in L^1$ and
$k_{\rm far}\in L^2\cap L^\infty$. For any nonnegative bounded
integrable spatial density $z$ with this mass and cap,
\begin{align}
 \|v[z]\|_\infty&\leq U_d\|k_{\rm near}\|_1
                         +M\|k_{\rm far}\|_\infty,\nonumber\\
 \|z^av[z]\|_1&\leq U_d^aM\|k_{\rm near}\|_1
                         +M^{3/2}\|k_{\rm far}\|_2.\label{eq:global-flux}
\end{align}
The near estimate uses $z^a\leq U_d^a$ and the $L^1$ convolution bound.
For the far estimate, Young and Cauchy--Schwarz give
$\|k_{\rm far}*z\|_2\leq M\|k_{\rm far}\|_2$ and
$\|z^a\|_2=\sqrt M$, since $a=1/2$.
Thus the source flux has a uniform spatial $L^1$ bound at almost every
time. Every compact smooth spatial pairing is absolutely continuous,
with its prescribed initial value.

Choose the continuous representatives simultaneously on a countable
dense family of compact smooth tests on increasing balls. Positivity
and the cap extend from almost every time to all times by continuity.
At each time these pairings define a density $z_*(t)$ between zero and
$U_d$, with $z_*(0)=d$. Nonnegative compact cutoffs show
$\int z_*(t)\leq M$ at every time. It agrees with the original solution
at almost every time. We now prove the tightness needed for equality of
mass and continuity against arbitrary bounded tests.

Let $\varphi(x)=\sqrt{1+|x|^2}-1$, so $|\nabla\varphi|\leq1$.
Choose smooth nondecreasing concave caps $\vartheta_L$ of the identity,
with derivative between zero and one, constant $c_L$ at infinity, and
$\vartheta_L\uparrow\mathrm{id}$. They can be obtained by scaling one
concave cap; its ratio to its argument is decreasing, which gives the
monotonicity in $L$. The function
$\vartheta_L(\varphi)-c_L$ is a compact smooth test. At good times add
$c_L$ times the known mass and use \eqref{eq:global-flux} to obtain
\[
 \int z(t)\vartheta_L(\varphi)\leq\int d\varphi+Kt.
\]
At an arbitrary time the compact part is continuous and the constant
part is at most $c_LM$, by the preliminary upper mass bound. Hence the
same upper bound holds for $z_*(t)$. Monotone convergence gives
\begin{equation}\label{eq:moment}
 \int z_*(t,x)\varphi(x)\,\dd x\leq\int d\varphi+Kt,
 \qquad t\geq0.
\end{equation}
This yields uniform tails on every bounded time interval. Approach an
arbitrary time by good times, pass a compact cutoff using continuity,
and let the cutoff radius increase using \eqref{eq:moment}. The result
is $\int z_*(t)=M$ at every time. For $\psi\in L^\infty$, first choose
a ball with uniformly small tails, then approximate $\psi$ there in
$L^1$ by a compact smooth function with bounded supremum. The cap
controls the on-ball error and \eqref{eq:moment} the exterior error.
The smooth pairing is continuous, so every bounded-test pairing is
continuous, including at zero. This representative is unique among
bounded-test continuous representatives of the same space-time solution.

The convex power norm is weak-$L^1$ lower semicontinuous, by its dual
variational formula with bounded tests. The pressure energy is lower
semicontinuous as well: each Fourier value is a bounded sine/cosine
pairing, and Fatou applies to the nonnegative integral
$\tfrac12\int|\xi|^{-2s}|\widehat z_*(\xi)|^2\dd\xi$.
Approach a terminal $t>0$ by increasing good times. The nonnegative
dissipation integrals increase to their full values. Both source
inequalities consequently hold at every terminal time. A countable
dense set of exponents and Fatou in the positive pair form give all
real $q>1$ simultaneously if needed. The equation and flux are unchanged
by the representative choice. These are all clauses of the stated class.
\end{proof}

\begin{proof}[Proof of Theorem~\ref{thm:main}]
The fixed $\alpha$ chosen in \eqref{eq:strict-margin} produces, by
Sections~\ref{sec:construction}--\ref{sec:shock}, a solution $u_c$ on
$[0,T_1]$ with the exact mass, cap, bounded-test continuity, pressure
inequality and every-real-$q$ power inequality. It vanishes outside
$B_{R(t)}$ at every such time, and $R$ is finite, positive and $C^1$,
with $R(0)=R_0$. Its final slice $d=u_c(T_1)$ is nonzero, nonnegative,
bounded and compactly supported; its pressure energy is finite.
Proposition~\ref{prop:global-source} supplies a global continuation
$z_*$ from this $d$, with the same parameters. Set
\[
 u(t)=\begin{cases}u_c(t),&0\leq t\leq T_1,\\
 z_*(t-T_1),&t>T_1.\end{cases}
\]
Split a compact space-time test at $T_1$. The equal terminal and initial
pairings with $d$ cancel, by their bounded-test continuity. This gives
the whole-space weak equation from $F$, with locally integrable flux.
All bounded pairings and the mass match at the junction, and the cap
remains $U$.

For $t>T_1$ add the pressure inequality for $u_c$ at $T_1$ to the one
for $z_*$ at $t-T_1$. The common energy $\mathcal E_\alpha(d)$ cancels,
giving the required inequality from $F$ and the full integral over
$(0,t)$. For each $q>1$ the identical coefficient
$4q(q-1)/(a+q)^2$ and common value $\|d\|_q^q$ give the same
cancellation in the power inequality. For $t\leq T_1$ the inequalities
have already been proved. This uses only inequalities from each
segment's own initial time.

Take $T_*=T_1$. The support assertion is required only for $0<t<T_*$.
At, for example,
$t=T_*/2$, take the finite radius $R(t)$. Nonnegativity and compact
support give $\int_{|x|>R(t)}u(t,x)\dd x=0$.
Since $s=1-\alpha/2\in(7/8,1)$ and $m=3/2$, this supplies one
admissible global solution contradicting the universal assertion.
\end{proof}


\section*{AI-use statement}
AI tools, including OpenAI Codex, were used in developing the argument, checking cited source
applications and preparing the manuscript. The paper is presented under
the author's responsibility.

\bibliographystyle{plainnat}
\bibliography{references}
\end{document}
